\documentclass[a4paper, 12pt]{article}

\usepackage[utf8]{inputenc}
\usepackage[T1]{fontenc}
\usepackage{textcomp}
\usepackage{newtxtext} \usepackage{newtxmath}
\usepackage{parskip}

\begin{document}

\begin{itemize}
    \item Adin Ballou 
    \item Petr Chelčický
\end{itemize}

Curiously, I have read more exegetical books about the Bible than books of the
Bible itself. I grew up in an atheist home, where the Bible was considered
nothing more than a matter of cultural importance. I devoted myself to Saint
Augustine, Origen of Alexandria, or Pseudo-Dionysius the Areopagite, and was
very fond of making and sharing ideas with religious friends, all at least a
decade before I even considered reading the Scriptures. And even when I finally
decided to read the Bible, I read only the Old Testament.

Now I am finally reading the New Testament, inspired by Tolstoi's work
\textit{The Kingdom of God is Within You}. As my readers know---\textit{scilicet
nulli}---I embrace the principles of libertarian socialism, and I came across
this book perchance when reading about Christian anarchism. Reading it, it
quickly became apparent that it made no sense to do so without examining the
Gospels myself. And so I looked for my King James Bible and opened the
\textit{Book of Matthew}.

Never would I have expected the reading to be so exquisitely beautiful, nor so
transformative, humbling, and illuminating. So much was the impact of this
perfect piece of literature, and of the perfect ethical doctrine which is
thereby exposed by Jesus Christ, that---however embarrassing it may be to
acknowledge it---I often found myself shedding tears as my eyes struggled to
move from one word onto another. And I, who have so often been guilty of pride
and arrogance, and who always found myself so secured in the vanity of
intelligence or knowledge, felt that I was the most ignorant and arrogant of
men. For I was finding now a radical and overcoming ethics---one in perfect
accordance with my own libertarian principles, but far more overreaching---from
the teachings of a poor carpenter, a man whose ethics vastly surpassed mine in
depth and breadth, and which more perfectly exposed the virtues of tolerance and
kindness.

It is a shame that the teachings of this extraordinary man have come to us
cloaked in rigid dogmatism and warped by the inherent evil which pervades every
human institution and all authority. And it is impossible for me not to observe,
having received the teachings of Christ at face value, and not through the
dogmatic interpretation of any institution, how shamelessly unchristian most
self-professed Christian people and institutions are. Surely their only claim
for Christianity can proceed from their belief in the supernatural elements of
the Bible---e.g. their belief in reincarnation or creationism---rather than
their practice of an ethical doctrine. 

But what is this doctrine? What are teachings of Jesus, as taught by Jesus
himself and not this or that church? In what do they differ from the values
which I already upheld, if in anything at all, and in what regards are they more
complete? 

Perhaps the most important principle which is derived from Jesus' teachings, and
one that is certainly not a necessary corollary of the values that sustain
libertarian socialism, is non-resistance to evil. I have embraced this principle
because I firmly believe it is an all-comprehensive and perfect doctrine. Its
concept and philosophy are well exposed in a long passage of Adin Ballou's work,
which I quote (with omissions) from \textit{The kingdom of God is within you}:

\begin{quote}
    Q. Whence originated the term “non-resistance?”

    A. From the injunction, “Resist not evil,” Matt. v. 39.

    Q. What does the term signify?

    A. It expresses a high Christian virtue, prescribed by Christ.

    Q. Is the word “resistance” to be taken in its widest meaning, that is, as showing that
    no resistance whatever is to be shown to evil? 

    A. No, it is to be taken in the strict sense of the Savior’s injunction; that is, we are
    not to retaliate evil with evil. Evil is to be resisted by all just means, but never with
    evil.

    Q. From what can we see that Christ in such cases prescribed non-resistance? 

    A. From the words which He then used. He said, “Ye have heard that it hath been
    said, An eye for an eye, and a tooth for a tooth. But I say unto you that ye resist not
    evil; but whosoever shall smite thee on thy right cheek, turn to him the other also.
    And if any man will sue thee at the law, and take away thy coat, let him have thy
    cloak also.” 


    Q. Did the ancients authorize the resistance of insult with insult?

    A. Yes; but Jesus prohibited this. A Christian has under no condition the right to
    deprive of life or to subject to insult him who does evil to his neighbor.

    Q. May a man kill or maim another in self-defence?

    A. No.

    Q. May he enter a court with a complaint, to have his insulter punished? 

    A. No; for what he is doing through others, he is in reality doing in his own
    person. 

    Q. May he fight with an army against enemies, or against domestic rebels? 

    A. Of course not. He cannot take any part in war or warlike preparations. He cannot
    use death-dealing arms. He cannot resist injury with injury, no matter whether he
    be alone or with others, through himself or through others. 

    Q. May he choose or fit out military men for the government? 

    A. He can do nothing of the kind, if he wishes to be true to Christ’s law. 

    Q. May he voluntarily give money, to aid the government, which is supported by
    military forces, capital punishment, and violence in general? 

    A. No, if the money is not intended for some special object, just in itself, where the
    aim and means are good. 

    Q. May he pay taxes to such a government? 

    A. No; he must not voluntarily pay the taxes, but he must also not resist their
    collection. The taxes imposed by the government are collected independently of
    the will of the subjects. It is impossible to resist the collection, without
    having recourse to violence; but a Christian must not use violence, and so he
    must give up his property to the violence which is exerted by the powers. 

    Q. May a Christian vote at elections and take part in a court or in the
    government? 

    A. No; the participation in elections, in the court, or in the government, is a
    participation in governmental violence. 

    Q. In what does the chief significance of the doctrine of non-resistance
    consist? 

    A. In that it alone makes it possible to tear the evil out by the root, both out of
    one’s own heart and out of the neighbor’s heart. This doctrine forbids doing that by
    which evil is perpetuated and multiplied. He who attacks another and insults him,
    engenders in another the sentiment of hatred, the root of all evil. To offend another,
    because he offended us, for the specious reason of removing an evil, means to repeat
    an evil deed, both against him and against ourselves — to beget, or at least to free, to
    encourage, the very demon whom we claim we wish to expel. Satan cannot be driven
    out by Satan, untruth cannot be cleansed by untruth, and evil cannot be vanquished
    by evil.

    True non-resistance is the one true resistance to evil. It kills and finally destroys the
    evil sentiment. 

    Q. But, if the idea of the doctrine is right, is it practicable? 

    A. It is as practicable as any good prescribed by the Law of God. The good cannot
    under all circumstances be executed without self-renunciation, privation, suffering,
    and, in extreme cases, without the loss of life itself. But he who values life more
    than the fulfilment of God’s will is already dead to the one true life. Such a man, in
    trying to save his life, shall lose it. Besides, in general, where non-resistance costs
    the sacrifice of one life, or the sacrifice of some essential good of life,
    resistance costs thousands of such sacrifices. Non-resistance preserves,
    resistance destroys. 

    It is incomparably safer to act justly than unjustly; to bear an insult than to
    resist it with violence — it is safer even in relation to the present life. If
    all men did not resist evil with evil, the world would be blessed. 

    Q. But if only a few shall act thus, what will become of them? 

    A. If only one man acted thus, and all the others agreed to crucify him, would
    it not be more glorious for him to die in the triumph of non-resisting love,
    praying for his enemies, than to live wearing the crown of Cæsar, bespattered
    with the blood of the slain? But one or thousands who have firmly determined not
    to resist evil with evil, whether among the enlightened or among savage
    neighbors, are much safer from violence than those who rely on violence. A
    robber, murderer, deceiver, will more quickly leave them alone than those who
    resist with weapons. They who take the sword perish with the sword, and those
    who seek peace, who act in a friendly manner, inoffensively, who forget and
    forgive offences, for the most part enjoy peace or, if they die, die blessed.

    Thus, if all kept the commandment of non-resistance, it is evident that there would be
    no offences, no evil deeds. If these formed a majority, they would establish the reign
    of love and good-will, even toward the ill-disposed, by never resisting evil with evil,
    never using violence. If there were a considerable minority of these, they would have
    such a corrective, moral effect upon society that every cruel punishment would be
    abolished, and violence and enmity would be changed to peace and love. If there were
    but a small minority of them, they would rarely experience anything worse than the
    contempt of the world, and the world would in the meantime, without noticing it,
    and without feeling itself under obligation, become wiser and better from this secret
    influence. And if, in the very worst case, a few members of the minority should be
    persecuted to death, these men, dying for the truth, would leave behind them their
    teaching, which is already sanctified by their martyr’s death.

    Peace be with all who seek peace, and all-conquering love be the imperishable
    inheritance of every soul, which voluntarily submits to the Law of Christ:
    Resist not evil.

\end{quote}

I do not wish here to convince others of this doctrine of which I am myself
convinced. But I should wish briefly to continue considering the teachings of
Jesus, insofar as his ethical doctrine derives principles for social action. For
example, why not recall the eschatological passage from \textit{Matthew}
25:35-40, where Jesus speaks of his return:

\begin{quote}
    Then shall the King say unto them on his right hand, Come, ye blessed of my
    Father, inherit the kingdom prepared for you from the foundation of the world: 
    For I was an hungred, and ye gave me meat: I was thirsty, and ye gave me
    drink: I was a stranger, and ye took me in: Naked, and ye clothed me: I was
    sick, and ye visited me: I was in prison, and ye came unto me. Then shall the
    righteous answer him, saying, Lord, when saw we thee an hungred, and fed thee?
    or thirsty, and gave thee drink? When saw we thee a stranger, and took thee in? or naked, and clothed thee? Or when saw we thee sick, or in prison, and came unto thee?
    And the King shall answer and say unto them, Verily I say unto you, Inasmuch as
    ye have done it unto one of the least of these my brethren, ye have done it unto
    me.
\end{quote}

Or is it not said in \textit{Ephesians}, 6:12, that 

\begin{quote}
    we wrestle not against flesh and blood, but against principalities, against
    powers, against the rulers of the darkness of this world, against spiritual
    wickedness in high places.
\end{quote}

And are not power and wealth the devil's offer, which Christ rejected, as
accounted in \textit{Luke} 4:6:

\begin{quote}
    And the devil said unto him, All this power will I give thee, and the glory
    of them: for that is delivered unto me; and to whomsoever I will I give it. 
    ---Luke 4:6
\end{quote}

or in \textit{Matthew} 4:1-11: 

\begin{quote}
    Again, the devil taketh him up into an exceeding high mountain, and sheweth him
    all the kingdoms of the world, and the glory of them; And saith unto him, All
    these things will I give thee, if thou wilt fall down and worship me. Then saith
    Jesus unto him, Get thee hence, Satan: for it is written, Thou shalt worship the
    Lord thy God, and him only shalt thou serve.
\end{quote}

But let me pause now: I do not wish to dwell too much in the Scriptures, however
many the passages which support my thesis---which is, I should say, not
original. I have become convinced that an honest reading of the New Testament,
one not pervaded by institutional dogma, leads to an absolute accordance with
the libertarian and altruistic principles which are the basis of socialist
thought---in a broad sense---. Here I wish to expose, not to prove: for
whosoever should wish a proof of this needs only consult the words of Jesus with
unprejudiced eyes. «He that hath ears to hear, let him hear». 

So what are the principles of the ethical doctrine which I uphold, the one I
have upheld so far and which has now been enriched and extended by the teachings
of Jesus? 

\begin{enumerate}
    \item Non-resistance to evil.
    \item Opposition to all terrestrial authority and concentrated power as inherently evil.
    \item Voluntary austerity, or rather a voluntary abstention from the pursuit
        of wealth, power, or sumptuosity.
    \item An abstention from judging others, as taught in Matthew 7:1-5, and
        from all arrogance and vanity.
    \item To forgive everything, as taught in Matthew 18:21-22.
    \item To side with those that are hungered, unclothed, imprisoned, subdued,
        and weak, and to go unto them.
    \item To never, under any condition, allow oneself to limit the liberty of a
        fellow man or woman; be it economic, political, religious, sexual, etc.
    \item To never worry: not about one's own life nor what tomorrow might
        bring, as taught in Matthew 6:25-34.
\end{enumerate}

Except for (1), I had already considered and favored all these principles
before, though I should say with less profoundness than now. It is (4) the one
that I fail the most to uphold, being so shamefully ill-disposed to judgment.



























\end{document}



